\section{同时对角化}
定理：对于每一个厄米算符$\Omega$，(至少)存在一个由其正交本征矢组成的基。它在这个本征基下是对角的，并且它的本征值就作为它的对角元。\par
定理：如果$\Omega ,\Lambda $是两个对易的厄米算符，则（至少）存在一个共同的本征态能把它们二者同时对角化\par
一个小补充：所谓对易，用物理语言来说就是可以同时观测到\par
在这里我们应该就可以看出来了，如果我们选定了由算符$\Omega$生成的基，如果另一个算符$\Lambda $是与其对易的，那么我们就可以得到一个本征态的两个信息量(请注意，我们始终在原来的空间中，即基不变)\par
既然我们可以得到一个本征态的两个信息量，很自然地，我们就可以对这个本征态进行标记(如角动量，用轨道角动量 \textit{l} 和磁量子数 \textit{m} 来表示一个态)